\documentclass[11pt,a4paper]{article}
\usepackage[utf8]{inputenc}
\usepackage{amsmath}
\usepackage{amsfonts}
\usepackage{amssymb}
\usepackage{graphicx}
\usepackage{booktabs}
\usepackage{hyperref}
\usepackage{geometry}
\geometry{margin=1in}

\title{Optimizing Question Generation via Discourse-Aware Salience and Hybrid Feature Fusion}
\author{Research Group - SQuAD-RST Project}
\date{\today}

\begin{document}

\maketitle

\begin{abstract}
This report details the implementation and evaluation of a Question Generation (QG) pipeline that utilizes Rhetorical Structure Theory (RST), linguistic features, and LLM-as-a-Judge ranking to identify salient sentences. We compare three modes: Zero-Shot, Hybrid Salient (RoBERTa + RST), and LLM-as-a-Judge. Our findings indicate that while LLMs often disagree with human-labeled "gold" sentences in SQuAD, the questions generated from LLM-selected salient sentences achieve higher qualitative scores and linguistic richness.
\end{abstract}

\section{Introduction}
Question Generation (QG) often suffers from "over-selection," where models generate too many redundant or trivial questions. We propose a "Salience First" architecture that filters paragraphs into the top-K most important sentences before generation.

\section{Methodology}
We evaluate three strategies for salience detection:
\begin{itemize}
    \item \textbf{Zero-Shot}: Direct question generation from the full paragraph without explicit salience filtering.
    \item \textbf{Hybrid Salient}: A RoBERTa-based classifier augmented with RST (Span Importance) and linguistic features (Surprisal, Complexity).
    \item \textbf{LLM Salient}: A paragraph-level ranking strategy using \texttt{Phi-4-mini-instruct} to select the top-2 most answerable sentences.
\end{itemize}

\section{Experimental Results}
The following table summarizes the performance across 10 sample paragraphs (retained 6 after cleaning) using BERTScore, ROUGE-L, and METEOR.

\begin{table}[h!]
\centering
\begin{tabular}{@{}lcccc@{}}
\toprule
\textbf{Mode} & \textbf{ROUGE-L} & \textbf{BERTScore-F1} & \textbf{METEOR} & \textbf{LLM Clarity} \\ \midrule
Zero-Shot      & 0.098            & 0.875                 & 0.131           & 4.50                 \\
Hybrid Salient & 0.176            & \textbf{0.883}        & 0.202           & 4.67                 \\
LLM Salient    & \textbf{0.176}   & 0.878                 & \textbf{0.233}  & \textbf{0.483}       \\ \bottomrule
\end{tabular}
\caption{Quantitative evaluation of generated questions per mode.}
\end{table}

\section{Discussion}
\subsection{The Gold Label Mismatch}
In our evaluation, the LLM Salient mode achieved 0\% recall against the SQuAD gold labels. However, it achieved the highest qualitative scores in METEOR and Clarity. This suggests that LLMs prioritize \textit{thematic importance}, whereas SQuAD gold labels often reflect \textit{arbitrary factual details} selected by human crowdworkers.

\subsection{Hybrid Efficiency}
The Hybrid Salient model provided the best balance between speed and performance, matching the LLM in ROUGE-L while being significantly more efficient to run.

\section{What's Next?}
\begin{enumerate}
    \item \textbf{Scale Up}: Run the full pipeline on 100+ samples to confirm the statistical significance of the METEOR/Clarity gains.
    \item \textbf{Few-Shot Tuning}: Incorporate SQuAD-specific examples into the LLM Ranking prompt to align its "salience" definition with human-labeled factual detail.
    \item \textbf{Recursive RST}: Explore using the hierarchy of the RST tree (Nucleus vs. Satellite) to weight the generation probability.
\end{enumerate}

\end{document}
